\section{Gothic Christianity}

John Ruskin, (pp. 189-191), {Bible of Amiens\footnote{\url{http://www.gutenberg.org/files/24428/24428-h/24428-h.htm}}}:

\begin{quotationx}
Quite the most beautiful sign of the power of true Christian-Catholic faith is this continual acknowledgement by it of the brotherhood -- nay, more, the fatherhood, of the elder nations who had not seen Christ; but had been filled with the Spirit of God; and obeyed, according to their knowledge, His unwritten law. The pure charity and humility of this temper are seen in all Christian art, according to its strength and purity of race; but best, to the full, seen and interpreted by the three great Christian-Heathen poets, Dante, Douglas of Dunkeld, and George Chapman. The prayer with which the last ends his life's work is, so far as I know, the perfectest and deepest expression of Natural Religion given us in literature; and if you can, pray it here -- standing on the spot where the builder once wrote the history of the Parthenon of Christianity:

\emph{\textbf{I Pray thee, Lord, the Father, and the Guide of our reason, that we may remember the nobleness with which Thou has adorned us; and that Thou wouldst be always on our right hand and on our left, in the motion of our own Wills: that so we may be purged from the contagion of the Body and the Affections of the Brute, and overcome them and rule; and use, as it becomes men to use them, for instruments. And then, that Thou would'st be in Fellowship with us for the careful correction of our reason, and for its conjunction by the light of truth with the things that truly are. And in the third place, I pray to Thee the Saviour, that thou would'st utterly cleanse away the closing gloom from the eyes of our souls, that we may know well who is to be held for God, and who for Mortal. Amen}.}
\end{quotationx}


John Ruskin wasn't much on Protestantism, although one senses that he could have acknowledged the innately Gothic, conservative nature that plays through Northern European spirituality:

\begin{quotex}
First, that throughout the Sermon on this Amiens Mount, Christ never appears, or is for a moment thought of, as the Crucified, nor as the Dead: but as the Incarnate Word---as the present Friend---as the Prince of Peace on Earth, and as the Everlasting King in Heaven. What His life \emph{is}, what His commands \emph{are}, and what His judgment \emph{will be}, are the things here taught: not what He once did, nor what He once suffered, but what He is now doing---and what He requires us to do. That is the pure, joyful, beautiful lesson of Christianity; and the fall from that faith, and all the corruptions of its abortive practice, may be summed briefly as the habitual contemplation of Christ's death instead of His Life, and the substitution of His past suffering for our present duty. \textbf{} Then, secondly, though Christ bears not \emph{His} cross, the mourning prophets,---the persecuted apostles---and the martyred disciples \emph{do} bear theirs. For just as it is well for you to remember what your undying Creator is \emph{doing} for you---it is well for you to remember what your dying fellow-creatures \emph{have done:} the Creator you may at your pleasure deny or defy---the Martyr you can only forget; deny, you cannot. Every stone of this building is cemented with his blood, and there is no furrow of its pillars that was not ploughed by his pain. \textbf{} Keeping, then, these things in your heart, look back now to the central statue of Christ, and hear His message with understanding. He holds the Book of the Eternal Law in His left hand; with His right He blesses,---but blesses on condition. ``This do, and thou shalt live''; nay, in stricter and more piercing sense, This \emph{be} and thou shalt live: to show Mercy is nothing---thy soul must be full of mercy; to be pure in act is nothing---thou shalt be pure in heart also. And with this further word of the unabolished law---''This if thou do \emph{not}, this if thou art not, thou shalt die.'' \textbf{} Die (whatever Death means)---totally and irrevocably. There is no word in thirteenth-century Theology of the pardon (in our modern sense) of sins; and there is none of the Purgatory of them. Above that image of Christ with us, our Friend, is set the image of Christ over us, our Judge. For this present life---here is His helpful Presence. After this life---there is His coming to take account of our deeds, and of our desires in them; and the parting asunder of the Obedient from the Disobedient, of the Loving from the Unkind, with no hope given to the last of recall or reconciliation. I do not know what commenting or softening doctrines were written in frightened minuscule by the Fathers, or hinted in hesitating whispers by the prelates of the early Church. But I know that the language of every graven stone and every glowing window,---of things daily seen and universally understood by the people, was absolutely and alone, this teaching of Moses from Sinai in the beginning, and of St. John from Patmos in the end, of the Revelation of God to Israel. This it was, simply---sternly---and continually, for the great three hundred years of Christianity in her strength (eleventh, twelfth, and thirteenth centuries), and over the whole breadth and depth of her dominion, from Iona to Cyrene,---and from Calpe to Jerusalem. At what time the doctrine of Purgatory was openly accepted by Catholic Doctors, I neither know nor care to know. It was first formalized by Dante, but never accepted for an instant by the sacred artist teachers of his time---or by those of any great school or time whatsoever. Neither do I know nor care to know---at what time the notion of Justification by Faith, in the modern sense, first got itself distinctively fixed in the minds of the heretical sects and schools of the North. Practically its strength was founded by its first authors on an asceticism which differed from monastic rule in being only able to destroy, never to build; and in endeavouring to force what severity it thought proper for itself on everybody else also; and so striving to make one artless, letterless, and merciless monastery of all the world. Its virulent effort broke down amidst furies of reactionary dissoluteness and disbelief, and remains now the basest of popular solders and plasters for every condition of broken law and bruised conscience which interest can provoke, or hypocrisy disguise. With the subsequent quarrels between the two great sects of the corrupted church, about prayers for the Dead, Indulgences to the Living, Papal supremacies, or Popular liberties, no man, woman, or child need trouble themselves in studying the history of Christianity: they are nothing but the squabbles of men, and laughter of fiends among its ruins. The Life, and Gospel, and Power of it, are all written in the mighty works of its true believers: in Normandy and Sicily, on river islets of France and in the river glens of England, on the rocks of Orvieto, and by the sands of Arno. But of all, the simplest, completest, and most authoritative in its lessons to the active mind of North Europe, is this on the foundation stones of Amiens. Believe it or not, reader, as you will: understand only how thoroughly it \emph{was} once believed; and that all beautiful things were made, and all brave deeds done in the strength of it---until what we may call `this present time,' in which it is gravely asked whether Religion has any effect on morals, \textbf{by persons who have essentially no idea whatever of the meaning of either Religion or Morality}. All human creatures, in all ages and places of the world, who have had warm affections, common sense, and self-command, have been, and are, Naturally Moral. Human nature in its fulness is necessarily Moral,---without Love, it is inhuman, without sense, inhuman,---without discipline, inhuman. In the exact proportion in which men are bred capable of these things, and are educated to love, to think, and to endure, they become noble,---live happily---die calmly: are remembered with perpetual honour by their race, and for the perpetual good of it. All wise men know and have known these things, since the form of man was separated from the dust. The knowledge and enforcement of them have nothing to do with religion: a good and wise man differs from a bad and idiotic one, simply as a good dog from a cur, and as any manner of dog from a wolf or a weasel. And if you are to believe in, or preach without half believing in, a spiritual world or law---only in the hope that whatever you do, or anybody else does, that is foolish or beastly, may be in them and by them mended and patched and pardoned and worked up again as good as new---the less you believe in---and most solemnly, the less you talk about---a spiritual world, the better. But if, loving well the creatures that are like yourself, you feel that you would love still more dearly, creatures better than yourself---were they revealed to you;---if striving with all your might to mend what is evil, near you and around, you would fain look for a day when some Judge of all the Earth shall wholly do right, and the little hills rejoice on every side; if, parting with the companions that have given you all the best joy you had on Earth, you desire ever to meet their eyes again and clasp their hands,---where eyes shall no more be dim, nor hands fail;---if, preparing yourselves to lie down beneath the grass in silence and loneliness, seeing no more beauty, and feeling no more gladness---you would care for the promise to you of a time when you should see God's light again, and know the things you have longed to know, and walk in the peace of everlasting Love---\emph{then}, the Hope of these things to you is religion, the Substance of them in your life is Faith. And in the power of them, it is promised us, that the kingdoms of this world shall yet become the kingdoms of our Lord and of His Christ.
\end{quotex}


John Ruskin here does much to delineate the spirit which Duby investigated\footnote{\url{https://www.amazon.com/Age-Cathedrals-Art-Society-980-1420/dp/0226167704/ref=la\_B001HD1HRE\_1\_6?ie=UTF8\&qid=1353721991\&sr=1-6}} in the early cathedral-building, the spirit which (if one accepts the claim being made here) represented the summit of Christianity, right after the dawn age of the Frankish imperium.

I am unaware of any rivals to this claim that are traditional in their orientation, excepting the Orthodox perspective which we dealt with last week. The modern critique of ``the Dark Ages'' hangs entirely on whether or not one accepts the modern technique of infinite manipulation (including of ``man's nature'', which is not supposed to exist, but yet is subject to intense discussion and engineering anyway). The Victorian era itself did not pretend to rival this period, but rather saw itself in its shadow, or else rejected it entirely. The Reformation period either borrowed heavily from it and lived in its clothes, or else attempted to purify it (or obliterate it). When we arrive at the last century, it was supposed to have disappeared forever, and the American era helped to deliberately deepen this delusion.

But no one, even the detractors, really believes that any following Western period can compete with it in terms of youth and springtime power, for it is my distinct impression that not even the Renaissance conceived of itself entirely as a new beginning (as the Dark Age was) but rather as a ``re-birth'' of a life that was already in existence.

What do the readers think?


\flrightit{Posted on 2012-11-23 by Logres}

\begin{center}* * *\end{center}

\begin{footnotesize}
\begin{sffamily}
\texttt{Jason-Adam on 2012-11-25 at 00:05 said:}

Ruskin's description of the unconscious collective mind of the Moyen Age is a perfect representation of the springtime of the West in the Spenglerian sense.

Two facts are made apparent : Frankish Christianity, not anything else, is the basis of Western, Franco-Germano-Roman culture, the culture of the restored Roman Empire that emerged from the ruins of the Vth and VIth centuries, and all further Western cultural life has been either an intellectualisation of or else a denial of this inner soul of the Christian West. The ``American'' period we know live in can be shown to be in fact a non-Western occupation of Western souls and minds... ... I am confident in the future of the West because history shows all cultures have a Second Religiousness phase, where the springtime beliefs reemerge and overcome materialism and rationalism. Who knows, perhaps Gornahoor is a beginning to this end ? Judging by comparisons to other cultures, the formation of the Western empire and second religiousness is not yet to come until sometime around 2080-2200 AD... .

One error Guenon and Evola might made about the West I think is to judge it as a culture in formation against those who were fully fulfilled... .looking into the future, if the West can overcome Americanism at some point the future seems brighter for Tradition... .

\hfill

\texttt{Logres on 2012-11-27 at 13:26 said:}

Jason-Adam, the letters of AKC that have been linked to on this site might be as helpful to you as they were to me (although I've not finished them). Is anyone else interested in reviving the Forum?

\hfill

\texttt{Jason-Adam on 2012-11-27 at 22:59 said:}

Please show me the link to the letters as regretfully I have not seen them before... ..sadly I have not studied AKC as much as I should have, something I will rectify soon.

The Forum seems a bad idea for me, as I do not desire to debate but only to ask questions and learn. Having been a member formerly of several metaphysical Internet discussion groups, the arguing I encountered there made me want to stay quiet.

\hfill

\texttt{Logres on 2012-11-28 at 06:21 said:}

There may be a better link, but here is the one I was using:

\url{http://www.scribd.com/doc/106308792/Selected-Letters-of-A-K-Coomaraswamy-PART-1}

Cologero first put this up.You may be right about the Forum, although ours seemed to operate differently.

\hfill

\texttt{Graham on 2012-11-28 at 17:04 said:}

There was debate, but no bickering. I agree that we should make an effort to revive it. The intention was always to conduct the more involved discussions there.

\hfill

\texttt{Logres on 2012-11-29 at 06:25 said:}

Perhaps we should agree on either a subject or topic to discuss?


\hfill

\end{sffamily}
\end{footnotesize}

